\section{Implementation Details}

\label{sec:implementation}
% ══════════════════════════════════════════════════════════════════════════════

\subsection{Software Stack}

HabitatNet is implemented in Python 3.11 using PyTorch 2.2 and PyTorch
Geometric 2.5. Spatial data processing uses \texttt{rasterio},
\texttt{geopandas}, and \texttt{h3-py} for hexagonal grid construction.
The corridor algorithm uses \texttt{networkx} for graph operations and
\texttt{scipy.sparse} for efficient Laplacian computations.

\subsection{Training Protocol}

\begin{itemize}
  \item \textbf{Optimizer}: AdamW with learning rate $3 \times 10^{-4}$,
    weight decay $10^{-4}$.
  \item \textbf{Schedule}: Cosine annealing with warm restarts over 200
    epochs.
  \item \textbf{Batch size}: 16 subgraphs of 10{,}000 nodes each,
    sampled using ClusterGCN \citep{chiang2019cluster}.
  \item \textbf{Graph layers}: $L = 6$ message-passing layers with 128
    hidden dimensions.
  \item \textbf{Regularization}: Dropout (0.2), edge dropout (0.1),
    and DropMessage \citep{fang2023dropmessage}.
  \item \textbf{Loss weights}: $\lambda_1 = 0.1$, $\lambda_2 = 0.05$
    (selected by validation).
\end{itemize}

\subsection{Computational Requirements}

Training a continental-scale model requires approximately 48 GPU-hours on
an NVIDIA A100 (80\,GB). Inference for a 1-million-cell region takes
approximately 4 minutes. The corridor algorithm runs in $O(B \cdot N
\cdot M^2)$ time, where $B$ is the budget, $N$ is the number of cells,
and $M$ is the number of patches.


% ══════════════════════════════════════════════════════════════════════════════
